% Presentation/content.tex
% Основной контент презентации к защите диссертации

\section{Методология}

\begin{frame}{Формальная модель: атрибутивный метаграф}
  \begin{columns}[T]
    \column{0.55\textwidth}
    \begin{itemize}
      \item Вершины = сущности (процессы, файлы, сетевые соединения).
      \item Рёбра = отношения с атрибутами: время, PID, хэш, IP и др.
      \item Подграфы = реализации тактик MITRE ATT\&CK (TA0001–TA0040).
      \item Модель позволяет учитывать \textbf{контекст} и \textbf{последовательность} действий.
    \end{itemize}
    \column{0.45\textwidth}
    \centering
    \begin{tikzpicture}[node distance=1.6cm and 2cm, >=Stealth, thick,
        proc/.style={rectangle, draw=blue!70!black, fill=blue!10, rounded corners, minimum width=1.6cm, minimum height=0.7cm},
        file/.style={ellipse, draw=green!70!black, fill=green!10, minimum width=1.5cm, minimum height=0.6cm},
        net/.style={diamond, draw=red!70!black, fill=red!10, minimum size=0.8cm},
        edge/.style={->}
    ]
    \node[proc] (p1) {explorer.exe};
    \node[file, below left=of p1] (f1) {doc.pdf};
    \node[proc, below right=of p1] (p2) {powershell};
    \node[net, below=of p2] (n1) {C2};

    \draw[edge] (p1) -- node[above left, font=\scriptsize] {открыл} (f1);
    \draw[edge] (p1) -- node[above right, font=\scriptsize] {запустил} (p2);
    \draw[edge] (p2) -- node[right, font=\scriptsize] {подкл.} (n1);
    \end{tikzpicture}
    \vspace{-0.2cm}
    \tiny Пример: тактика «Execution» (TA0002)
  \end{columns}
\end{frame}

\begin{tikzpicture}[>=Stealth, font=\scriptsize]
  \definecolor{ta0002}{HTML}{99CC99}
  \definecolor{ta0003}{HTML}{9999FF}
  \definecolor{ta0005}{HTML}{FFCC66}
  \definecolor{ta0010}{HTML}{CCCC99}

  \node[fill=ta0002, circle, inner sep=3pt, label=above:{T1059}] (t2) at (0,0) {};
  \node[fill=ta0003, circle, inner sep=3pt, label=above:{T1078}] (t3) at (2,0) {};
  \node[fill=ta0005, circle, inner sep=3pt, label=below:{T1027}] (t4) at (1,-1.2) {};
  \node[fill=ta0010, circle, inner sep=3pt, label=below:{T1071.001}] (t5) at (3,-1.2) {};

  \draw[->, thick] (t2) -- (t3);
  \draw[->, thick] (t2) -- (t4);
  \draw[->, thick] (t3) -- (t5);

  % Легенда без matrix
  \node[font=\tiny, anchor=west] at (-2.5,-2) {TA0002: Execution};
  \node[font=\tiny, anchor=west] at (-2.5,-2.3) {TA0003: Persistence};
  \node[font=\tiny, anchor=west] at (-2.5,-2.6) {TA0005: Defense Evasion};
  \node[font=\tiny, anchor=west] at (-2.5,-2.9) {TA0010: Exfiltration};

  % Цветные квадратики
  \fill[ta0002] (-3.2,-2) rectangle (-2.9,-1.7);
  \fill[ta0003] (-3.2,-2.3) rectangle (-2.9,-2.0);
  \fill[ta0005] (-3.2,-2.6) rectangle (-2.9,-2.3);
  \fill[ta0010] (-3.2,-2.9) rectangle (-2.9,-2.6);
\end{tikzpicture}


\begin{comment}
\begin{frame}{Интеграция с MITRE ATT\&CK}
  \centering
  \begin{tikzpicture}[>=Stealth, font=\scriptsize]
    \definecolor{ta0002}{HTML}{99CC99}
    \definecolor{ta0003}{HTML}{9999FF}
    \definecolor{ta0005}{HTML}{FFCC66}
    \definecolor{ta0010}{HTML}{CCCC99}

    \node[fill=ta0002, circle, inner sep=3pt, label=above:{T1059}] (t2) at (0,0) {};
    \node[fill=ta0003, circle, inner sep=3pt, label=above:{T1078}] (t3) at (2,0) {};
    \node[fill=ta0005, circle, inner sep=3pt, label=below:{T1027}] (t4) at (1,-1.2) {};
    \node[fill=ta0010, circle, inner sep=3pt, label=below:{T1071.001}] (t5) at (3,-1.2) {};

    \draw[->, thick] (t2) -- (t3);
    \draw[->, thick] (t2) -- (t4);
    \draw[->, thick] (t3) -- (t5);

    \matrix [draw, below=1.8cm of t4, column sep=0.5em, row sep=0.2em, font=\tiny] {
      \node[fill=ta0002, circle, inner sep=2pt] {}; & \node{TA0002: Execution}; \\
      \node[fill=ta0003, circle, inner sep=2pt] {}; & \node{TA0003: Persistence}; \\
      \node[fill=ta0005, circle, inner sep=2pt] {}; & \node{TA0005: Defense Evasion}; \\
      \node[fill=ta0010, circle, inner sep=2pt] {}; & \node{TA0010: Exfiltration}; \\
    };
  \end{tikzpicture}
\end{frame}
\end{comment}

\section{Эксперименты}

\begin{frame}{Наборы данных}
    \begin{itemize}
        \item \textbf{MITRE APT29 (2020)} — целевые атаки на Windows/Linux.
        \item \textbf{Финансовые атаки} — логи от банковских партнёров (2023–2024).
        \item \textbf{Нормальный трафик} — корпоративная сеть MIREA (анонимизирована).
    \end{itemize}
  \vspace{0.5cm}
  \centering
  \begin{tabular}{lrr}
    \toprule
    Источник & События & Атаки \\
    \midrule
    APT29 & 1.2 млн & 18 \\
    Финансовые & 3.7 млн & 42 \\
    Нормальный трафик & 5.1 млн & 0 \\
    \bottomrule
  \end{tabular}
\end{frame}

\begin{frame}{Сравнение методов обнаружения}
  \centering
  \begin{tikzpicture}
    \begin{axis}[
        ybar,
        bar width=14pt,
        width=0.9\textwidth,
        height=5cm,
        symbolic x coords={Snort, Isolation Forest, GNN, Предложенный},
        xtick=data,
        ylabel={F1-score},
        ymin=0, ymax=1,
        nodes near coords,
        every node near coord/.append style={font=\scriptsize},
        axis lines=left,
        enlarge x limits=0.5,
        xticklabel style={font=\scriptsize, align=center, text width=2cm}
    ]
    \addplot[fill=gray!50] coordinates {
        (Snort, 0.49)
        (Isolation Forest, 0.64)
        (GNN, 0.75)
        (Предложенный, 0.87)
    };
    \end{axis}
  \end{tikzpicture}
\end{frame}

\begin{frame}{Производительность}
    \begin{itemize}
        \item Время обработки 1 млн событий: \textbf{8.2 с} (Intel i7-12700, 16 ГБ RAM).
        \item Поддержка потоковой обработки с окном 5 минут.
        \item Потребление памяти: \textbf{< 1.2 ГБ} даже при полной загрузке.
    \end{itemize}
  \vspace{0.5cm}
  \centering
  \footnotesize Результаты получены на реальных данных без синтетических упрощений.
\end{frame}

\begin{comment}
\section{Списки}
\begin{frame}[plain, noframenumbering]
    \begin{center}
        \Huge
        Списки
    \end{center}
\end{frame}

\subsection{Нумерованные}

\begin{frame}
    \frametitle{Нумерованные списки}
    \begin{enumerate}
        \item один
        \item два
        \item три
    \end{enumerate}
\end{frame}
\note{
    Этот текст будет виден только если его отображение включено
    в~файле \textbf{Presentation/setup}.
    Для раздельного вывода презентации и заметок на~разные экраны (как
    в~impress или powerpoint) можно использовать программу
    \textit{pdf-presenter-console}.
}

\subsection{Не нумерованные}


\begin{frame}
    \frametitle{Перечисления}
    \begin{itemize}
        \item Проблема 1
        \item Проблема 2
        \item Проблема 3
    \end{itemize}
\end{frame}
\note[itemize]{
    \item Тезис 1
    \item Тезис 2
    \item Тезис 3
}

\subsection{Комбинированные}

\begin{frame}
    \frametitle{Комбинация списков}
    \begin{enumerate}
        \item \textbf{Задача 1}
              \begin{itemize}
                  \item Подзадача 1-1
                  \item Подзадача 1-2
              \end{itemize}
        \item \textbf{Задача 2}
              \begin{itemize}
                  \item Подзадача 2-1
                  \item Подзадача 2-2
                  \item Подзадача 2-3
              \end{itemize}
        \item \textbf{Задача 3}
              \begin{itemize}
                  \item Подзадача 3-1
                  \item Подзадача 3-2
                  \item Подзадача 3-3
              \end{itemize}
    \end{enumerate}
\end{frame}
\note[itemize]{
    \item Задача 1
    \item Задача 2
    \item Задача 3
}

\begin{frame}[allowframebreaks]
    \frametitle{Разделение слайда}
    Поясняющий текст
    \begin{itemize}
        \item Один
        \item Два
        \item Три
    \end{itemize}
    \framebreak
    Продолжение предыдущего слайда
\end{frame}

\section{Графика}
\begin{frame}[plain, noframenumbering]
    \begin{center}
        \Huge
        Графика
    \end{center}
\end{frame}


\begin{frame}
    \frametitle{Одиночное изображение}
    \centering
    \includegraphics[width=0.8\linewidth]{latex} % окружение figure не требуется
\end{frame}

\begin{frame}
    \frametitle{Векторная графика}
    \begin{figure}
        \centering
        \ifdefmacro{\tikzsetnextfilename}{\tikzsetnextfilename{tikz_presentation}}{}% присваиваемое предкомпилированному pdf имя файла (не обязательно)
        \input{Presentation/images/tikz_plot.tikz}
    \end{figure}
\end{frame}

\subsection{Расположение}

\begin{frame}
    \frametitle{Изображения по-вертикали}
    \centering
    \vfill
    \includegraphics[width=0.8\linewidth,height=0.1\textheight]{latex} \\
    \TeX
    \vfill
    \includegraphics[width=0.8\linewidth,height=0.2\textheight]{latex} \\
    \LaTeX
    \vfill
    \includegraphics[scale=0.2]{latex} \\
    \vfill
\end{frame}


\begin{frame}
    \frametitle{Изображения по-горизонтали}
    \begin{minipage}[t]{0.47\linewidth}
        \textbf{Составная \\ подпись 1}
        \center{\includegraphics[width=1\linewidth]{knuth1}}
    \end{minipage}
    \hfill
    \begin{minipage}[t]{0.47\linewidth}
        \textbf{Составная \\ подпись 2}
        \center{\includegraphics[width=1\linewidth]{knuth2}}
    \end{minipage}
\end{frame}

\subsection{Линии}

\begin{frame}
    \frametitle{Разделяющие линии}
    \begin{minipage}[c]{0.47\linewidth}
        \center{\includegraphics[width=1\linewidth]{latex}}
        \bigskip
        \hrule{}
        \bigskip
        \textbf{Составная \\ подпись 1}
    \end{minipage}
    \hfill
    \vrule{}
    \hfill
    \begin{minipage}[c]{0.47\linewidth}
        \flushright
        \textbf{Составная \\ подпись 2}
        \center{\includegraphics[width=1\linewidth]{knuth2}}
    \end{minipage}
\end{frame}

\section{Остальное}
\begin{frame}[plain, noframenumbering]
    \begin{center}
        \Huge
        Остальное
    \end{center}
\end{frame}

\subsection{Формулы}

\begin{frame}
    \frametitle{Формулы}
    \[
    \left\{
    \begin{array}{rl}
        \dot x = & \sigma (y-x)  \\
        \dot y = & x (r - z) - y \\
        \dot z = & xy - bz
    \end{array}
    \right.
    \]
\end{frame}

\begin{frame}
    \frametitle{amsmath}
    \centering
    \begin{minipage}[t]{0.5\linewidth}
        \begin{multline*}
            y = 1 x^1 + 2 x^2 + 3 x^3 + \\ + 4 x^4 + 5 x^5 + \dots
        \end{multline*}
    \end{minipage}
\end{frame}

\begin{frame}[allowframebreaks]
    \frametitle{Уравнения Максвелла}
    \centering{
        \small
        \def\arraystretch{1.8}%
        \begin{tabular}{ll}
            \toprule
            Интегральная форма                                                                                                                                            & Дифференциальная форма                                                          \\ \midrule
            \(Q_e(t) = \displaystyle\oiint_S \vec D(t) \cdot d\vec{s} = \displaystyle\iiint_V \rho_v(t) dv\)                                                              & \(\nabla \cdot \vec D(t) = \rho_v(t)\)                                          \\
            \(\displaystyle\oiint_S \vec B(t) \cdot d\vec{s} = 0\)                                                                                                        & \(\nabla \cdot \vec B(t) = 0\)                                                  \\
            \(V_{emf}(t) = \displaystyle\oint_L \vec E(t) \cdot d\vec{l}\) = \(- \displaystyle\iint_S \left[\frac{\partial\vec{B}(t)}{\partial t}\right] \cdot d\vec{s}\) & \(\nabla \times \vec E(t) = - \frac{\partial\vec{B}(t)}{\partial t}\)           \\
            \(I(t) = \displaystyle\oint_L \vec H(t) \cdot d\vec{l} = \displaystyle\iint_S \left[\vec J(t) + \frac{\partial\vec{D}(t)}{\partial t}\right] \cdot d\vec{s}\) & \(\nabla \times \vec H(t) = \vec J(t) + \frac{\partial\vec{D}(t)}{\partial t}\) \\ \midrule
            \(\displaystyle\oiint_S \vec J \cdot d\vec{s} = -\frac{\partial Q_e}{\partial t}\)                                                                            & \(\nabla \cdot \vec J = - \frac{\partial \rho_v}{\partial t}\)                  \\
            \bottomrule
            \multicolumn{2}{c}{\(\vec D(t) = \left[\varepsilon(t)\right] * \vec E(t)\)}                                                                                                                                                                     \\
            \multicolumn{2}{c}{\(\vec B(t) = \left[\mu(t)\right] * \vec H(t)\)}                                                                                                                                                                             \\
        \end{tabular}
    }
    \framebreak
    
    \hspace{0.05\linewidth}
    \centering{
        \small
        \def\arraystretch{1.8}%
        \begin{tabular}{ll}
            \toprule
            Интегральная форма                                                                                                                & Дифференциальная форма                               \\ \midrule
            \(Q_e = \displaystyle\oiint_S \vec D \cdot d\vec{s} = \displaystyle\iiint_V \rho_v dv\)                                           & \(\nabla \cdot \vec D = \rho_v\)                     \\
            \(\displaystyle\oiint_S \vec B \cdot d\vec{s} = 0\)                                                                               & \(\nabla \cdot \vec B = 0\)                          \\
            \(V_{emf} = \displaystyle\oint_L \vec E \cdot d\vec{l}\) = \(- \displaystyle\iint_S \left[j \omega \vec B\right] \cdot d\vec{s}\) & \(\nabla \times \vec E = - j \omega \vec B\)         \\
            \(I = \displaystyle\oint_L \vec H \cdot d\vec{l} = \displaystyle\iint_S \left[\vec J + j \omega \vec D\right] \cdot d\vec{s}\)    & \(\nabla \times \vec H = \vec J + j \omega \vec{D}\) \\ \midrule
            \(\displaystyle\oiint_S \vec J \cdot d\vec{s} = - j \omega Q_e\)                                                                  & \(\nabla \cdot \vec J = - j \omega \rho_v\)          \\
            \bottomrule
            \multicolumn{2}{c}{\(\vec D(t) = \left[\varepsilon\right] \vec E(t)\)}                                                                                                                   \\
            \multicolumn{2}{c}{\(\vec B(t) = \left[\mu\right] \vec H(t)\)}                                                                                                                           \\
        \end{tabular}
    }
\end{frame}

\subsection{Таблицы}

\begin{frame}
    \frametitle{Таблица}
    \centering
    \begin{tabular}{|l|l|}
        \hline
        \textbf{Заголовок 1} & \textbf{Заголовок 2} \\
        \hline
        Сумма                & \(b+a\)              \\
        \hline
        Разность             & \(a-b\)              \\
        \hline
        Произведение         & \(a*b\)              \\
        \hline
    \end{tabular}
\end{frame}

\begin{frame}
    \frametitle{Другая таблица}
    \centering
    \begin{tabular}{lc}
        \toprule
        \multicolumn{1}{c}{\textbf{Заголовок 1}} & \textbf{Заголовок 2} \\ \midrule
        Сумма                                    & \(b+a\)              \\
        Разность                                 & \(a-b\)              \\
        Произведение                             & \(a*b\)              \\
        \bottomrule
    \end{tabular}
\end{frame}


\subsection{Разное}

\begin{frame}
    \frametitle{Большой многоуровневый список}
    \begin{itemize}
        \item \textbf{Пункт 1}
              \begin{itemize}
                  \itemi Подпункт 1-1
                  \itemi Подпункт 1-2
              \end{itemize}
        \item \textbf{Пункт 2}
              \begin{itemize}
                  \itemi Подпункт 2-1
              \end{itemize}
        \item \textbf{Пункт 3}
              \begin{itemize}
                  \itemi Подпункт 3-1
                  \itemi Подпункт 3-2
              \end{itemize}
        \item \textbf{Пункт 4}
              \begin{itemize}
                  \itemi Подпункт 4-1
              \end{itemize}
        \item \textbf{Пункт 5}
              \begin{itemize}
                  \itemi Подпункт 5-1
                  \itemi Подпункт 5-2
                  \itemi Подпункт 5-3
              \end{itemize}
    \end{itemize}
\end{frame}

\begin{frame}
    \frametitle{Четыре изображения}
    \centering
    \includegraphics[width=0.35\linewidth,angle=35]{latex}
    \includegraphics[width=0.35\linewidth,angle=135]{latex}\\
    \includegraphics[width=0.35\linewidth,angle=15]{latex}
    \includegraphics[width=0.35\linewidth,angle=-15]{latex}
\end{frame}
\end{comment}